\documentstyle[%


righttag%


,11pt%


,amssymb%


,russian%               remove in Eng.


,izvru%                 Set izven% in Eng.


]{amsart}





% 





\newcommand{\Arg}{\operatorname{Arg}}


\newcommand{\tts}[1]{}





\theoremstyle{definition}


\theorembodyfont{\rm}


\newtheorem{notenonum}{\hskip\parindent Замечание}


\newtheorem{alg}{\hskip\parindent Алгоритм}


\renewcommand{\thenotenonum}{}





%\hoffset=-15mm%                  For Eng.


\setcounter{page}{3}


\god{2002}


\nomer{12~(487)} %                        Remove in Eng.


%\nomer{Vol.\,46, No.\,12}%               Set in Eng.


%\str{\pageref{begin}--\pageref{end}}%  Set in Eng.


\udk{519.853}





\author{\it А.А.\,Андрианова, Я.И.\,Заботин}


% NB:   ^^ remove in Eng.


\title{Управление процессом минимизации \\


в параметризованном методе центров}





\date{}


%\pagestyle{myheadings}%         Set in Eng.


\pagestyle{plain} %              Remove in Eng.


\begin{document}


%\thispagestyle{plain}%          Set in Eng.


\maketitle


\label{begin}








Получение решения задачи математического программирования с заданной точностью


за конечное число итерационных шагов с практической точки зрения является 


актуальной вычислительной проблемой. В разработке алгоритмов, позволяющих 


решить эту проблему в методе центров [1], [2], условно можно выделить 


два основных подхода. Первый подход заключается в разработке алгоритмов, 


использующих специальные критерии остановки, условия которых выполняются 


через конечное число итераций и гарантируют достижение заданной


точности. В рамках этого подхода разработаны алгоритмы с двусторонним 


приближением к минимуму [3], с неполной минимизацией функции максимума 


[4], с аппроксимацией допустимого множества [5], [6].





Исследования в рамках второго подхода заключались во введении управляющих


параметров, позволяющих получить требуемую точность за заранее


определенное число итераций. Теоретически в [7]--[9] было показано, 


что параметризация вспомогательной функции максимума может позволить


найти приближенное с заданной точностью решение за один процесс


минимизации вспомогательной функции. 





Данная статья посвящена разработке алгоритмов в параметризованном методе 


центров с аппроксимацией множества допустимых решений в рамках второго 


подхода.





Пусть $H=\{1,2,\dotsc,m\}$, функции $f(x)$, $f_i(x)$, $i\in H$, определены в 


$n$-мерном евклидовом пространстве $R_n$, $g(x)=\max\{f_i(x),\ i\in H\}$, 


$D=\{x: x\in R_n,\  g(x)\le 0\}$, $D'=\{x: x\in R_n$, $g(x)<0\}$.


\medskip





{\bf Условие A.} Задача


\begin{equation}


\min\{f(x),\ x\in D\}


\end{equation}


удовлетворяет условию A, если она имеет решение, функции $f(x)$,


$g(x)$ непрерывны в $R_n$, $D'\ne\emptyset$ и $\overline{D'}=D$. 


\medskip





Введем следующие обозначения: $f^*=\min\{f(x),\ x\in D\}$,


$Y=\Arg\min\{f(x),\ x\in R_n\}$, $X_{\sigma}^*=\{x: x\in D,\ f(x)\le


f^*+\sigma\}$, где $\sigma\ge 0$. Очевидно, $X_0^*$ является множеством


точных решений задачи (1). 





Tребуется по заданному $\varepsilon>0$ 


найти точку из множества $X_{\varepsilon}^*$.





\begin{lem}


Если задача $(1)$ удовлетворяет условию $\mathrm A$ и для


заданной константы $\sigma>0$ множество $X_{\sigma}^*$ является


ограниченным, то $\min\{g(x),\ x\in X_{\sigma}^*\}<0$.


\end{lem}





\begin{pf}


Пусть $x^*\in X_0^*$. Так как $f(x^*)<f^*+\sigma$ при 


$\sigma>0$, то в силу непрерывности функции $f(x)$ найдется такая 


окрестность $\omega$ точки $x^*$, что $f(x)<f^*+\sigma$ для всех 


$x\in \omega$. Тогда $\omega\cap D\subset X_{\sigma}^*$. Кроме того, в


силу условия $\overline{D'}=D$ в $\omega$ найдется точка $y\in D'$.


Из двух последних предложений следует утверждение 


$y\in \omega\cap D'\subset \omega\cap D\subset X_{\sigma}^*$ и $g(y)<0$. 


\end{pf}





Для фиксированных значений параметров $t, \gamma, \rho$ положим


\begin{align*}


F(x,t,\gamma,\rho)&=\max\{f(x)-t,\ \rho g(x)-\gamma\},\\


%


Z(t,\gamma,\rho)&=\Arg\min\{F(x,t,\gamma,\rho),\ x\in R_n\}. 


\end{align*}





\begin{thm}


Пусть функции $f(x)$, $g(x)$ определены и непрерывны в


$R_n$, для константы $\varepsilon>0$ множество $X_{\varepsilon}^*$ является 


ограниченным, параметры $\rho>0$, $\gamma$, $t$


зафиксированы таким образом, что


\begin{equation}


\rho\min\{g(x),\ x\in X_{\varepsilon}^*\}-\varepsilon\le


f^*+\gamma-t\le 0.


\end{equation}


Тогда $Z(t,\gamma,\rho)\subset X_{\varepsilon}^*$.


\end{thm}





\begin{pf}


Зафиксируем произвольно точку $x^*\in X_0^*$.


Тогда имеем, во-первых, $\rho g(x^*)-\gamma\le -\gamma$, во-вторых,


согласно условиям теоремы $f(x^*)-t\le-\gamma$. Поэтому


$F(x^*,t,\gamma,\rho)=\max\{f(x^*)-t$, $\rho g(x^*)-\gamma\}\le-\gamma$.


Отсюда для любой точки $z\in Z(t,\gamma,\rho)$ имеют место неравенства


$\rho g(z)-\gamma\le F(z,t,\gamma,\rho)\le F(x^*,t,\gamma,\rho)\le


-\gamma$, откуда $g(z)\le 0$. Таким образом, $Z(t,\gamma,\rho)\subset D$.


Поэтому для доказательства теоремы достаточно показать, что 


$Q\cap Z(t,\gamma,\rho)=\emptyset$, где $Q=D\setminus X_{\varepsilon}^*$.





Согласно определению множества $X_{\varepsilon}^*$ для любой точки $y\in


X_{\varepsilon}^*$ выполняется неравенство $f(y)\le f^*+\varepsilon$, а


для любой точки $x\in Q$ имеет место $f(x)>f^*+\varepsilon$. Поэтому


справедливо двойное неравенство


\begin{equation}


f(x)-t>f^*+\varepsilon -t\ge f(y)-t \quad \forall x\in Q,\ \ \forall y\in


X_{\varepsilon}^*.


\end{equation}





Множество $X_{\varepsilon}^*$ является замкнутым по определению и


ограниченным согласно условиям теоремы. Непрерывная функция $g(x)$ 


достигает на нем своего минимального значения, т.\,е.


существует точка $\ov y \in X_{\varepsilon}^*$ такая, что $g(\ov


y)=\min\{g(x),\ x\in X_{\varepsilon}^*\}$. В силу (2) имеем $f^*+\gamma-t\ge


\rho\min\{g(x),\ x\in X_{\varepsilon}^*\}-\varepsilon$, поэтому для точки


$\ov y$ согласно первому соотношению из (3) получим цепочку неравенств 


$f(x)-t>f^*+\varepsilon-t\ge


\rho\min\{g(x),\ x\in X_{\varepsilon}^*\}-\gamma=\rho g(\ov y)-\gamma$ 


для всех $x\in Q$. Отсюда в силу (3) и из определения функции максимума


имеем 


\begin{equation}


F(x,t,\gamma,\rho)\ge f(x)-t>\max\{f(\ov y)-t,\ \rho g(\ov


y)-\gamma\}=F(\ov y,t,\gamma,\rho) \quad \forall x\in Q.


\end{equation}





Предположим теперь, что существует точка $z\in Z(t,\gamma,\rho)\cap Q$.


Тогда согласно (4) для этой точки выполняется неравенство


$F(z,t,\gamma,\rho)>F(\ov


y,t,\gamma,\rho)$, которое противоречит тому, что $z\in Z(t,\gamma,\rho)$. 


Поэтому $Z(t,\gamma,\rho)\cap Q=\emptyset$, что, как


отмечалось, достаточно для доказательства теоремы.


\end{pf}





Из условий теоремы 1 естественным образом вытекают оценки параметра $\rho$,


позволяющие при различной фиксации значений параметра $t$ и определенных


правилах фиксации значений параметра $\gamma$


получить решение с точностью $\varepsilon>0$ 


путем минимизации функции типа $F(x,t,\gamma,\rho)$.





\begin{cor}


Пусть задача (1) удовлетворяет условию A,


$\varepsilon>0$, $X_{\varepsilon}^*$ --- ограниченное множество и параметры


$t$, $\gamma$ зафиксированы так, что $\gamma\ge\delta>0$, $t\ge f^*+\gamma$.


Если параметр $\rho>0$ зафиксировать таким образом, что


\begin{equation}


\rho\ge\frac{\varepsilon+\delta+f^*-t}{\min\{g(x),\ x\in


X_{\varepsilon}^*\}},


\end{equation}


то $Z(t,\gamma,\rho)\subset X_{\varepsilon}^*$.


\end{cor}





\begin{pf}


Согласно лемме 1 выполняется неравенство $\min\{g(x),\ x\in


X_{\varepsilon}^*\}<0$. Поэтому из (5) следует $\rho\min\{g(x),\ x\in


X_{\varepsilon}^*\}\le \varepsilon+\delta+f^*-t$. Отсюда, т.\,к. по условиям


следствия $\gamma\ge\delta$, получим неравенство 


$\gamma\ge\delta\ge t-f^*-\varepsilon+\rho\min\{g(x)$,


$x\in X_{\varepsilon}^*\}$. 


Кроме того, согласно условиям следствия выполняется $f^*+\gamma-t\le0$.


Следовательно, имеет место двойное неравенство (2). По теореме 1 при


такой фиксации параметров $Z(t,\gamma,\rho)\subset X_{\varepsilon}^*$.


\end{pf}





\begin{notenonum}


Пусть имеет место случай, когда числитель в правой части дроби из


(5) неотрицательный, т.\,е. $t\le f^*+\delta+\varepsilon$. Это возможно,


например, когда значения $\gamma$, $\delta$ фиксируются так, что


$\delta\le\gamma\le\delta+\varepsilon$ и $t=f^*+\gamma$. Тогда дробь в правой


части (5) согласно лемме 1 оказывается неположительной и,


следовательно, $Z(t,\gamma,\rho)\subset X_{\varepsilon}^*$ при любом $\rho>0$. 


\end{notenonum}





\begin{cor}


Если задача (1) удовлетворяет условию A,


$\varepsilon>0$, $X_{\varepsilon}^*$ --- ограниченное множество 


и параметры $\overline{\rho}>0$, $t$, $\gamma$ зафиксированы так, что 


\begin{equation}


t\le f^*,\quad \overline{\rho}\ge\frac{\varepsilon-f^*+t}{\min\{g(x),\ x\in


X_{\varepsilon}^*\}}, \quad \gamma=\overline{\rho}\min\{g(x),\ x\in


X_{\varepsilon}^*\}-\varepsilon,


\end{equation}


то $Z(t,\gamma,\rho)\subset X_{\varepsilon}^*$ для любого


$\rho\ge\overline{\rho}$.


\end{cor}





\begin{pf}


В силу леммы 1 имеет место неравенство $\min\{g(x)$,


$x\in X_{\varepsilon}^*\}<0$. Отсюда и из второго неравенства (6) получим


$\overline{\rho}\min\{g(x),\ x\in X_{\varepsilon}^*\}\le\varepsilon-f^*+t$, 


что в силу третьего из соотношений (6) означает


\begin{equation}


f^*+\gamma-t\le 0.


\end{equation}


Так как $\rho\ge\ov\rho$, то равенство из (6) приводит к неравенству


$\gamma\ge\rho\min\{g(x),\ x\in X^*_\varepsilon\}-\varepsilon$,


а учитывая первое условие из (6), получим


$f^*+\gamma-t\ge\rho\min\{g(x)$, $x\in X_{\varepsilon}^*\}-\varepsilon$,


что вместе с (7) означает выполнение условий теоремы 1. Следовательно, 


$Z(t,\gamma,\rho)\subset X_{\varepsilon}^*$.


\end{pf}





\begin{notenonum}


Возможен случай, когда числитель правой части дроби из


второго неравенства (6) неотрицателен. Для этого достаточно зафиксировать


параметр $t$ так, чтобы выполнялось 


двойное неравенство $f^*-\varepsilon\le t\le f^*$. Тогда сама дробь окажется


неположительной, и $\varepsilon$-решение задачи (1) можно получить


как точку абсолютного минимума функции $F(x,t,\gamma,\rho)$ при любом


$\rho>0$ и параметре $\gamma$, зафиксированном согласно (6).


\end{notenonum}





Следствия 1 и 2 теоремы 1 указывают способы фиксации управляющих 


параметров для получения $\varepsilon$-решения при $t\ge f^*+\gamma$ и 


$t\le f^*$ соответственно. Следующие положения указывают, как получить 


$\varepsilon$-решение при фиксации значений $t$ и $\gamma$ такиx, что 


$f^*<t<f^*+\gamma$.





\begin{lem}


Пусть функции $f(x)$, $g(x)$ непрерывны в $R_n$. Если


$Y\cap D\ne\emptyset$ и


параметры $t$, $\gamma$ зафиксированы так, что $\gamma\ge t-f^*$, то 


$X_0^*\subset Z(t,\gamma,\rho)$ при любом $\rho>0$, причем


$Z(t,\gamma,\rho)\subset Y$.


\end{lem}





\begin{pf}


Пусть зафиксирован параметр $\rho>0$.


Выберем точку $x^*\in X_0^*$.


Поскольку $Y\cap D\ne\emptyset$, имеет место включение $X_0^*\subset Y$.


Tак как $x^*\in D$ и $\rho>0$, справедливо неравенство $\rho


g(x^*)\le 0$. Поэтому согласно условиям леммы справедливо


$f(x^*)-t\ge-\gamma\ge \rho g(x^*)-\gamma$. Следовательно,


$F(x^*,t,\gamma,\rho)=f(x^*)-t$. Так как $x^*\in Y$, для любых точек


$x\in R_n$ выполняется $f(x^*)\le f(x)$. Отсюда в силу определения


функции максимума получим следующую цепочку неравенств


$F(x^*,t,\gamma,\rho)=f(x^*)-t\le f(x)-t\le F(x,t,\gamma,\rho)$ для любых


$x\in R_n$. Следовательно, $x^*\in Z(t,\gamma,\rho)$. Поскольку точка


$x^*\in X_0^*$ была зафиксирована произвольным образом, первая часть


утверждения леммы доказана.





Докажем второе утверждение.


Для любой точки $z\in Z(t,\gamma,\rho)$ имеет место 


соотношение $f(z)-t\le


F(z,t,\gamma,\rho)\le F(x^*,t,\gamma,\rho)=f(x^*)-t$, а т.\,к. $x^*\in


Y$, все неравенства в последней цепочке выполняются как равенства, т.\,е.


$f(z)=f^*$. Таким образом, $Z(t,\gamma,\rho)\subset Y$.


\end{pf}





Лемма 2 показывает, что если среди допустимых решений задачи (1)


есть точки абсолютного минимума функции $f(x)$ и параметры $t$, $\gamma$


зафиксированы так, что $t\le f^*+\gamma$, то при любом положительном


значении $\rho$ среди точек абсолютного минимума функции


$F(x,t,\gamma,\rho)$ существует та, которая является точным решением


задачи (1).





Обозначим через $M$ множество функций, определенных в $R_n$, каждый


локальный минимум которых является абсолютным.





\begin{lem}


Пусть функции $f(x)$, $g(x)$ непрерывны в $R_n$, $f\in M$,


$\rho>0$. Если $Y\cap D=\emptyset$


и параметры $t$, $\gamma$ зафиксированы так, что


$\gamma>t-f^*$, то $Z(t,\gamma,\rho)\cap


D=\emptyset$.


\end{lem}





\begin{pf}


Предположим, что утверждение леммы не имеет места,


т.\,е. существует точка $z\in Z(t,\gamma,\rho)\cap D$. Тогда $g(z)\le


0$, $f(z)\ge f^*$ и согласно условиям леммы имеет место цепочка неравенств 


$f(z)-t\ge f^*-t>-\gamma\ge \rho g(z)-\gamma$. В силу непрерывности функций 


$f(x)$ и $g(x)$ существует окрестность $\omega$ точки $z$ такая, что 


$f(x)-t\ge \rho g(x)-\gamma$ для всех точек $x\in\omega$. Поэтому, т.\,к.


$z\in Z(t,\gamma,\rho)$, имеет место 


$f(z)-t\le F(z,t,\gamma,\rho)\le F(x,t,\gamma,\rho)=f(x)-t$


для всех $x\in \omega$. Следовательно, $z$ является точкой


локального минимума функции $f(x)$, а поскольку $f\in M$, то $z\in


Y$. Однако доказанное противоречит условию $Y\cap D=\emptyset$.


Полученное противоречие и завершает доказательство.


\end{pf}





\begin{notenonum}


Пусть в задаче (1) $f\in M$. Выберем произвольно точку 


$z\in Z(t,\gamma,\rho)$ при значениях параметров, зафиксированных так, что


$0<\gamma\le\varepsilon$, $f^*<t<f^*+\gamma$, $\rho>0$.


Пусть $z\notin D$, что справедливо в силу леммы 3, если $Y\cap D=\emptyset$, 


или в случае, когда $Y\cap D\ne\emptyset$, но $Y\setminus D\ne\emptyset$. 


Тогда в силу фиксации параметра $\gamma$


имеем $f^*<t<f^*+\gamma\le f^*+\varepsilon$.


Если же $z\in D$, то $Y\cap D\ne\emptyset$, поскольку в противном случае


согласно лемме 3 такую точку получить нельзя. Поэтому в силу леммы 2


точка $z$ является точным решением задачи (1). 


\end{notenonum}





Oценки (5) и (6) применять на практике непосредственно нельзя, т.\,к. они 


содержат неизвестные величины. Однако при наложении дополнительных условий на 


исходные данные задачи можно получить оценки неизвестных величин, 


используемых в (5) и (6), и тем самым получить более грубые, но 


практически реализуемые варианты этих оценок. 





Для получения оценки значения $f^*$ сверху достаточно, например, выбрать 


точку $\overline{z}\in D$. Тогда $f(\overline{z})\ge f^*$. 





Укажем способ построения оценки снизу величины $f^*$. Минимизируем


функцию $\Psi(x)=\max\{f(x)+d, g(x)\}$, где $d$ --- константа, выбранная


так, что $f(x)+d\ge g(x)$ для всех $x\in D$. Тогда согласно лемме 1


[3], если точка $y\in \Arg\min\{\Psi(x),\ x\in R_n\}$ такова, что


$y\notin D$, то $f(y)\le f^*$. В противном случае точка $y$ является 


точным решением задачи (1).





Числа, оценивающие значение $f^*$ сверху и снизу, будем обозначать


$\overline{f}$ и $\underline{f}$ соответственно.





Для построения оценки величины $\min\{g(x),\ x\in X_{\varepsilon}^*\}$


потребуется наложить дополнительные условия на данные исходной задачи и


доказать несколько вспомогательных утверждений.





\begin{lem}


Пусть функции $\varphi(x)$, $\psi(x)$ определены и


непрерывны в $R_n$, $\varphi\in M$, $G=\{x: x\in R_n,\ \psi(x)\le


0\}\ne\emptyset$, $Q=\Arg\min\{\varphi(x),\ x\in


G\}$. Если $\min\{\varphi(x),\ x\in G\}\ne\min\{\varphi(x),\ x\in R_n\}$,


то $\psi(z)=0$ для любых $z\in Q$.


\end{lem}





\begin{pf}


Если для любой точки $x\in G$ имеет место


$\psi(x)=0$, то утверждение леммы тривиально выполняется. Пусть


существует точка $x\in G$, для которой $\psi(x)<0$. Предположим,


что утверждение леммы неверно и существует точка $z\in Q$, для которой


$\psi(z)<0$. Тогда $G$ является такой окрестностью точки $z$, что


$\varphi(z)\le\varphi(x)$ для всех $x\in G$. Следовательно, $z$ является


точкой локального минимума функции $\varphi(x)$, а т.\,к. $\varphi \in


M$, то она является и точкой абсолютного минимума. Однако по условиям


леммы абсолютный и глобальный минимумы на множестве $G$ функции


$\varphi(x)$ не совпадают.


\end{pf}





\begin{lem}


Пусть задача $(1)$ удовлетворяет условию $\mathrm A$, $g\in M$,


$G_1=\{x: x\in R_n$, $g(x)\le 0$, $g_1(x)\le 0\}\ne\emptyset$, где


$g_1(x)$ --- непрерывная в $R_n$ функция, $Q_1=\Arg\min\{g(x),\ x\in G_1\}$. 


Если $\min\{g(x),\ x\in G_1\}\ne\min\{g(x),\ x\in R_n\}$, то 


$g_1(z)=0$ для любых $z\in Q_1$.


\end{lem}





\begin{pf}


Положим $\Phi(x)=\max\{ g(x), g_1(x)\}$. Тогда


$G_1=\{ x: x\in R_n$, $\Phi(x)\le 0\}$. Согласно лемме 4 \ $\Phi(z)=0$ для


всех $z\in Q_1$, т.\,е.


\begin{equation}


\max\{ g(z), g_1(z)\}=0.


\end{equation}





Если $g_1(z)=0$ для любой точки $z\in Q_1$, то утверждение леммы 


справедливо. Допустим, существует точка $z\in Q_1$ такая, что


\begin{equation}


g_1(z)<0.


\end{equation}


Тогда $g(z)=0$. Так как $\overline{D'}=D$, в окрестности


$\omega=\{x: x\in R_n,\ g_1(x)<0\}$ точки $z$ существует точка $y\in D'$.


Отсюда $g(y)<0$ и в силу $g_1(y)<0$ имеем $y\in G_1$. Так как $z\in


Q_1$, то $g(z)\le g(y)<0$, что вместе с (9) противоречит (8).


\end{pf}





{\bf Условие B.} Задача (1) удовлетворяет


условию A, функция $f(x)$ является выпуклой на


$R_n$, и функции $f_i(x)$, $i\in H$, являются сильно выпуклыми на $R_n$ с 


постоянными сильной выпуклости~$\mu_i$.





\begin{thm}


Пусть задача $(1)$ удовлетворяет условию $\mathrm B$,


$\varepsilon>0$, множество $X_{\varepsilon}^*$ ограничено, функция $f(x)$ 


удовлетворяет условию Липшица с константой $L$ на


множестве $X_{\varepsilon}^*$ и $\min\{g(x)$, $x\in


X_{\varepsilon}^*\}\ne\min\{g(x)$, $x\in R_n\}$.


Тогда имеет место неравенство


\begin{equation}


\min\{g(x), x\in X_{\varepsilon}^*\}\le-\frac{\mu\varepsilon^2}{L^2},


\end{equation}


где $\mu=\min\{\mu_i,\ i\in H\}$.


\end{thm}





\begin{pf}


По определению множество $X_{\varepsilon}^*$


является замкнутым, а по условиям теоремы --- ограниченным множеством.


Поэтому непрерывная функция $g(x)$ достигает на нем своего минимального


значения, т.\,е. существует точка $y\in X_{\varepsilon}^*$, для которой


$g(y)=\min\{g(x)$, $x\in X_{\varepsilon}^*\}$. Из определения сильной


выпуклости ([10], с.\,181) следует, что функция $g(x)$ является сильно 


выпуклой с постоянной сильной выпуклости $\mu=\min\{\mu_i,\ i\in


H\}$. Так как функция $f(x)$ выпуклая, множество $X_{\varepsilon}^*$ 


является выпуклым. Известно ([10], с.\,182), что в


этом случае имеет место неравенство $\mu\|x-y\|^2\le g(x)-g(y)$ для всех


$x\in X_{\varepsilon}^*$, в том числе и для точки $x^*\in X_0^*$,


т.\,к. согласно определению $X_0^*\subset X_{\varepsilon}^*$. Далее,


поскольку $g(x^*)\le 0$, получим неравенство $\mu\|x^*-y\|^2\le


-g(y)$ или


\begin{equation}


g(y)\le -\mu\|x^*-y\|^2.


\end{equation}


Так как $\overline{D'}=D$, то множество $X_{\varepsilon}^*$ содержит


точку из $D'$. Тогда $f(y)=f^*+\varepsilon$


согласно определению $X_{\varepsilon}^*$ и в силу 


леммы 5, а т.\,к. $f(x)$


удовлетворяет условию Липшица, имеет место 


$\varepsilon=f^*+\varepsilon-f^*=|f(y)-f(x^*)|\le L\|y-x^*\|$ или


$\|x^*-y\|\ge \frac{\varepsilon}{L}$. Отсюда и из (11) получим


$\min\{g(x)$, $x\in X_{\varepsilon}^*\}=g(y)\le -\mu\|x^*-y\|^2\le


-\frac{\mu\varepsilon^2}{L^2}$.


\end{pf}





Исходя из дополнительной информации о задаче, получим практически


реализуемые аналоги оценок (5) и (6).





\begin{thm}


Пусть задача $(1)$ удовлетворяет условию $\mathrm B$,


$\varepsilon>0$, множество 


$X_{\varepsilon}^*$ ограничено, функция $f(x)$ удовлетворяет на


$X_{\varepsilon}^*$ условию Липшица с константой $L$ и 


$\min\{g(x)$, $x\in X_{\varepsilon}^*\}\ne\min\{g(x)$, $x\in R_n\}$. Если


параметры $\rho>0$, $t$, $\gamma$ зафиксированы таким образом, что 


\begin{equation} 


\gamma\ge\delta>0, \quad t\ge f^*+\gamma, \quad


\rho\ge-\frac{L^2(\varepsilon+\delta+\underline{f}-t)}{\mu\varepsilon^2},


\end{equation}


где $\mu=\min\{\mu_i,\ i\in H\}$, то 


$Z(t,\gamma,\rho)\subset X_{\varepsilon}^*$. 


\end{thm}





\begin{pf}


Умножим третье неравенство (12) на величину


$-\frac{\mu\varepsilon^2}{L^2}<0$. Отсюда и из оценки $\underline{f}\le


f^*$ получим цепочку неравенств


$-\frac{\mu\varepsilon^2}{L^2}\rho\le\varepsilon+\delta+\underline{f}-t


\le\varepsilon+\delta+f^*-t$. Принимая к тому же во внимание (10) и


условие $\rho>0$, получим $\rho\min\{g(x)$, $x\in


X_{\varepsilon}^*\}\le-\rho\frac{\mu\varepsilon^2}{L^2}\le


\varepsilon+\delta+f^*-t$, что, в свою очередь, в силу леммы 1


приводит к неравенству (5), которое


вместе с соотношениями (12) согласно следствию


1 к теореме 1 означает, что


$Z(t,\gamma,\rho)\subset X_{\varepsilon}^*$.


\end{pf}





\begin{thm}


Пусть задача $(1)$ удовлетворяет условию $\mathrm B$,


$\varepsilon>0$, множество $X_{\varepsilon}^*$ ограничено, функция $f(x)$ 


удовлетворяет на $X_{\varepsilon}^*$ условию Липшица с константой $L$ 


и $\min\{g(x)$, $x\in X_{\varepsilon}^*\}\ne\min\{g(x)$, $x\in R_n\}$. Если


параметры $\rho'>0$, $t$, $\gamma$ зафиксированы таким образом, что 


\begin{equation} 


t\le f^*, \quad


\rho'\ge-\frac{L^2(\varepsilon-\overline{f}+t)}{\mu\varepsilon^2}, \quad


\gamma=-\rho'\frac{\mu\varepsilon^2}{L^2}-\varepsilon,


\end{equation}


где $\mu=\min\{\mu_i,\ i\in H\}$,


то $Z(t,\gamma,\rho)\subset X_{\varepsilon}^*$


для любых $\rho\ge\rho'$.


\end{thm}





\begin{pf}


Пусть


\begin{equation}


\eta=-\frac{\rho'\mu\varepsilon^2}{L^2\min\{g(x), x\in


X_{\varepsilon}^*\}}.


\end{equation}


В силу леммы 1 выполняется $\eta>0$. С другой


стороны, в силу (10) имеет место неравенство


$-\frac{\mu\varepsilon^2}{L^2\min\{g(x),


\ x\in X_{\varepsilon}^*\}}\le 1$, из которого, т.\,к. $\rho'>0$, следует


\begin{equation}


0<\eta\le\rho'.


\end{equation}


Из (14) с учетом второго из неравенств (13) получим соотношение


$$


-\frac{\eta L^2\min\{g(x),\ x\in


X_{\varepsilon}^*\}}{\mu\varepsilon^2}=\rho'\ge-\frac{L^2(\varepsilon-


\overline{f}+t)}{\mu\varepsilon^2},


$$


откуда $\eta\min\{g(x),\ x\in


X_{\varepsilon}^*\}\le \varepsilon-\overline{f}+t$. А так как


$\overline{f}\ge f^*$, то в силу леммы 1


\begin{equation}


\eta\ge\frac{\varepsilon-f^*+t}{\min\{g(x),\ x\in X_{\varepsilon}^*\}}.


\end{equation}


Подставляя $\rho'$ в третье из соотношений (13), получим


$\gamma=\eta\min\{g(x),\ x\in X_{\varepsilon}^*\}-\varepsilon$, что


согласно следствию 2 к теореме 1 вместе с (16) и первым неравенством из


(13) означает,


что $Z(t,\gamma,\rho)\subset X_{\varepsilon}^*$ для любых $\rho\ge\eta$,


в том числе согласно (15) для всех $\rho\ge\rho'$.


\end{pf}





Правила фиксации параметров, определенные теоремами 3 и 4, являются


практически реализуемыми и позволяют


получить $\varepsilon$-решение задачи (1) с помощью одного процесса


минимизации функции вида $F(x,t,\gamma,\rho)$. Однако значение параметра 


$\rho$, используемое при этом, может оказаться достаточно большим, что 


может привести к овражному характеру минимизируемой 


функции $F(x,t,\gamma,\rho)$ и трудностям при ее минимизации.


Следующие алгоритмы позволяют уменьшить влияние больших значений


параметра $\rho$ и получить $\varepsilon$-решение задачи (1) не более


чем за заданное заранее количество $N$ процессов минимизации


вспомогательной функции вида $F(x,t,\gamma,\rho)$. Каждый такой процесс


будем называть большой итерацией.





\begin{alg}


Зафиксируем произвольно точку $x_0\in D$, требуемую


точность $\varepsilon>0$, число $\delta\in(0;\varepsilon]$, $N>0$ ---


число больших итераций, за которое должно быть получено $\varepsilon$-решение


задачи (1), $\rho_{-1}=0$. Полагаем $k=0$.





1. Вычисляются значения


$C_k=-\frac{L^2(\varepsilon+\delta+\underline{f}-f(x_k))}{\mu\varepsilon


^2}$, $\rho_k=\max\{\frac{k+1}{N}C_k,\ \rho_{k-1}\}$. 





2. Фиксируется число $\gamma_k\in [\delta; \varepsilon]$.





3. Выбирается точка $x_{k+1}\in Z(f(x_k),\gamma_k,\rho_k)$.





4. Если $k=N-1$ и $x_{k+1}\in D$, то $x_{k+1}\in X_{\varepsilon}^*$.





5. Если $x_{k+1}\notin D$, то $x_k\in X_{\varepsilon}^*$, иначе переход


к п.\,1 при $k$, замененном на $k+1$.


\end{alg}





\begin{thm}


Пусть задача $(1)$ удовлетворяет условию $\mathrm B$, 


множество $X_{\varepsilon}^*$ ограничено, функция $f(x)$ удовлетворяет


на $X_{\varepsilon}^*$ условию Липшица с константой $L$ и 


$\min\{g(x)$, $x\in X_{\varepsilon}^*\}\ne\min\{g(x)$, $x\in R_n\}$,


$\mu=\min\{\mu_i$, $i\in H\}$. Тогда


условия п.\,$4$ или п.\,$5$ алгоритма $1$ выполнятся при некотором $k\le N-1$


и будет получено $\varepsilon$-решение задачи $(1)$.


\end{thm}





\begin{pf}


Пусть для некоторого номера $k\le N-1$


выполняются условия п.\,5 алгоритма 1, т.\,е. $x_{k+1}\notin D$. Тогда


$x_k\in D$, поскольку условия п.\,5 алгоритма


выполняются впервые для номера $k$. Поэтому


согласно утверждению 4 леммы 1 [5] и


$\gamma_k\le\varepsilon$ имеем $f(x_k)-f^*\le\gamma_k\le\varepsilon$. Отсюда


следует, что $x_k\in X_{\varepsilon}^*$.





Пусть теперь $x_{k+1}\in D$ для всех номеров $k\le N-1$. Тогда для


номера $k=N-1$ имеем $\rho_k=\max\{C_k,\ \rho_{k-1}\}\ge C_k$. Отсюда в


силу определения значения $C_k$ и согласно теореме 3 следует, что


$x_{k+1}\in X_{\varepsilon}^*$.


\end{pf}





\begin{notenonum}


Известно, что при решении задачи (1) методом внутренних


центров строится последовательность точек $\{x_k\}$, для которой для


любых $k$ имеет место $f(x_{k+1})\le f(x_k)$. Если рассматривать $C_k$


как функцию от $f(x_k)$, то такая функция является возрастающей.


Следовательно, $C_{k+1}\le C_k$ для любого номера $k$. Таким образом,


нет необходимости увеличивать значение параметра $\rho_k$ до значения


$C_0$, что позволяет на практике избежать трудностей, связанных с


овражным характером функции $F(x,f(x_k),\gamma_k,\rho_k)$. 


\end{notenonum}





\begin{alg}


Зафиксируем произвольно точку $x_0\notin D$, для


которой $f(x_0)\le f^*$, требуемую точность $\varepsilon>0$, 


число $\delta\in(0;\varepsilon]$, $N>0$ ---


число больших итераций, гарантирующее получение $\varepsilon$-решения


задачи (1), $\rho_{-1}=0$. Полагаем $k=0$.





1. Вычисляются значения


$C_k=-\frac{L^2(\varepsilon-\overline{f}+f(x_k))}{\mu\varepsilon


^2}$, $\rho_k=\max\{\frac{k+1}{N}C_k,\ \rho_{k-1}\}$. 





2. Фиксируется число


$\gamma_k=\varepsilon+\frac{\mu\varepsilon^2}{L^2}\rho_k$.





3. Выбирается точка $x_{k+1}\in Z(f(x_k),-\gamma_k,\rho_k)$.





4. Если $x_{k+1}\in D$, то $x_{k+1}\in X_{\varepsilon}^*$, иначе --- переход


к п.\,1 при $k$, замененном на $k+1$. 


\end{alg}





\begin{thm}


Пусть задача $(1)$ удовлетворяет условию $\mathrm B$, 


$Y\cap D=\emptyset$, множество $X_{\varepsilon}^*$ ограничено, функция


$f(x)$ удовлетворяет условию Липшица на $X_{\varepsilon}^*$ с константой


$L$ и 


$\min\{g(x)$, $x\in X_{\varepsilon}^*\}\ne\min\{g(x)$, $x\in R_n\}$,


$\mu=\min\{\mu_i,\ i\in H\}$. Тогда


при некотором $k\le N-1$ выполнятся условия п.\,$4$ алгоритма $2$ 


и будет получено $\varepsilon$-решение задачи $(1)$.


\end{thm}





\begin{pf}


Если для любого $k<N-1$ условия п.\,4 не


выполняются, то для номера $k=N-1$ имеет место


$\rho_k=\max\{C_k,\ \rho_{k-1}\}\ge C_k$, что в силу теоремы 4 означает,


что $x_{k+1}\in X_{\varepsilon}^*$, причем для этого номера выполняются


условия п.\,4 алгоритма 2, поскольку $X_{\varepsilon}^*\subset D$. 





Предположим теперь, что для некоторого номера $k<N-1$ имели место


условия п.\,4 алгоритма 2, т.\,е. $x_{k+1}\in D$. Докажем, что


$x_{k+1}\in X_{\varepsilon}^*$. 





Имеет место неравенство 


\begin{equation}


-\gamma_k\le f(x_k)-f^*.


\end{equation}


Действительно, допустим, что справедливо противное.


Тогда, т.\,к. $Y\cap D=\emptyset$, то


$$


Z(f(x_k),-\gamma_k,\rho_k)\cap D=\emptyset


$$


согласно лемме 3,


что невозможно в силу выбора точки $x_{k+1}\in


Z(f(x_k),-\gamma_k,\rho_k)$ такой, что $x_{k+1}\in D$.





Предположим, что $f^*-\gamma_k-f(x_k)<\rho_k\min\{g(x)$,


$x\in X_{\varepsilon}^*\}-\varepsilon$.


Тогда в силу (10) и правила фиксации параметра


$\gamma_k$ из п.\,2 алгоритма 2 имеем


$f^*-f(x_k)<\gamma_k-\varepsilon+\rho_k\min\{g(x)$, $x\in


X_{\varepsilon}^*\}\le\gamma_k-\varepsilon-\rho_k\frac{\mu\varepsilon^2}


{L^2}=0$. Таким образом, $f(x_k)>f^*$, но, с другой стороны, согласно


лемме 1 [5], т.\,к. $x_k\notin D$, имеем $f(x_k)\le f^*$. Получено


противоречие. Поэтому $f^*-\gamma_k-f(x_k)\ge


\rho_k\min\{g(x)$, $x\in X_{\varepsilon}^*\}-\varepsilon$.


Данное неравенство вместе с (17) в силу теоремы 1 


означает, что $x_{k+1}\in X_{\varepsilon}^*$.


\end{pf}





\begin{notenonum}


Известно, что при решении задачи (1) методом внешних


центров строится последовательность точек $\{x_k\}$, для которой для


любых $k$ имеет место $f(x_{k+1})\ge f(x_k)$.


Величину $C_k$ можно рассматривать


как функцию от аргумента $f(x_k)$. Эта функция является убывающей.


Следовательно, $C_{k+1}\le C_k$ для любого номера $k$. Поэтому, как и


в замечании к алгоритму 1, отметим, что можно избежать вычислительных


трудностей при минимизации функции $F(x,f(x_k),-\gamma_k,\rho_k)$,


поскольку нет необходимости увеличения значения параметра $\rho_k$ до 


$C_0$.


\end{notenonum}





За $N$ больших итераций, проделанных по алгоритмам 1 и 2, будет 


гарантировано получение


$\varepsilon$-решения задачи (1). Однако оба алгоритма имеют критерии


остановки (п.\,5 алгоритма 1 и п.\,4 алгоритма 2), позволяющие


прекратить вычислительный процесс на большой итерации, по номеру меньшей


$N$. Результаты проведенных вычислительных экспериментов показали, что 


остановка процесса вычислений по алгоритмам 1 и 2 чаще 


происходила по выполнению условий указанных критериев,


нежели по выполнению всех $N$ больших итераций. 





\begin{thebibliography}{99}





\bibitem{1}


Huard P. {\it Resolution of mathematical programming with


nonlinear constraints by the method of centres} // Non. Progr. --


1967. -- P.\,207--219.


\bibitem{2}


Гроссман К., Каплан А.А. {\it Нелинейное программирование на


основе безусловной минимизации}. -- Новосибирск: Наука, 1981. -- 183 с.


\bibitem{3}


Заботин Я.И., Даньшин И.Н. {\it Алгоритмы с комбинированием,


параметризацией и двусторонним приближением в методе центров} // Изв.


вузов. Математика. -- 1998. -- \No\,12. -- С.\,40--48.


\bibitem{4}


Заботин Я.И., Даньшин И.Н. {\it Алгоритмы в методе центров с неполной


минимизацией вспомогательной функции максимума} // Изв. вузов.


Математика. -- 1999. -- \No\,12. -- С.\,53--59.


\bibitem{5}


Заботин Я.И., Андрианова А.А. {\it Алгоритмы в методе центров с


аппроксимацией множества допустимых решений} // Изв. вузов. Математика.


-- 2001. -- \No\,12. -- С.\,41--45.


\bibitem{6}


Андрианова А.А. {\it Конечные алгоритмы в методе центров с


аппроксимацией допустимого множества}. -- Казанск. ун-т. -- Казань, 2002. --


31 с. -- Деп. в ВИНИТИ 06.05.02, \No\,791-B2002.


\bibitem{7}


Заботин Я.И. {\it Минимаксный метод решения задачи математического


программирования} // Изв. вузов. Математика. -- 1975. -- \No\,6.


-- С.\,36--43.


\bibitem{8}


Заботин Я.И., Князев Е.А. {\it Алгоритмы с адаптацией в


параметризованном методе центров} // Исследов. по прикл. матем. -- Казань,


1987. -- \No\,14. -- С.\,9--15. 


\bibitem{9}


Князев Е.А. {\it Алгоритмы с адаптацией в методе центров}: Дис.


$\dotsc$ канд. физ.-матем. наук. -- Казань, 1987. -- 134 с.


\bibitem{10}


Васильев Ф.П. {\it Численные методы решения экстремальных задач}. -- 


Москва: Наука, 1988. -- 552 с.





\end{thebibliography}





\vskip 2cm





\post{\it Казанский государственный}{\it Поступила}


\post{\it университет}{27.08.2002}





\label{end}


\end{document}


